\section*{Erd\H{o}s Problem 1140: finiteness of universally prime quadratic differences}

\paragraph{Statement and scope.}
There are only finitely many positive integers \(n\) for which \(n-2x^2\) is prime whenever \(x\) is a positive integer with \(2x^2<n\).
The following proof establishes finiteness, without claiming an explicit final list or resolving a possible exceptional integer in stronger classification results. See \url{https://www.erdosproblems.com/1140}. It is a reconstruction via the prime-character input used in the modern general resolution of the related problem 1141.

\paragraph{External prime-character theorem.}
Pollack's Theorem 1.3 in \emph{Bounds for the first several prime character nonresidues}, \url{https://arxiv.org/abs/1508.05035}, states that for every \(\varepsilon,A>0\) and all sufficiently large moduli \(M\), a quadratic Dirichlet character modulo \(M\) has at least \((\log M)^A\) primes \(p\leq M^{1/4+\varepsilon}\) with character value \(1\). We use its nonprincipal case, with \(\varepsilon=1/16\).

If \(2n\) is not a square, the Kronecker character
\(\chi(t)=(8n/t)\) is a nonprincipal quadratic Dirichlet character modulo \(M=8n\).
The theorem produces an odd prime \(p\nmid2n\), with
\[
 p\leq(8n)^{5/16},\qquad (2n/p)=1.
\]
Since multiplication by \(2^2\) does not change a nonzero quadratic-residue class, this implies that \(2x^2\equiv n\pmod p\) has a solution \(1\leq r<p\).
For sufficiently large \(n\),
\[
 2r^2\leq2p^2=O(n^{5/8})=o(n),
 \qquad n-2r^2>p.
\]
The positive integer \(n-2r^2\) is divisible by \(p\) and is greater than \(p\), so it is composite.

\paragraph{The principal-character case.}
If \(2n\) is a square, choose any odd prime \(p\nmid2n\). It can be chosen with \(p=n^{o(1)}\): among the first \(\omega(2n)+1\) odd primes at least one does not divide \(2n\), and repeated Bertrand postulates bound that prime by \(2^{O(\omega(2n))}\).
Here \(\omega(n)=o(\log n)\), since the product of \(r\) distinct primes is at least \((r+1)!\).
Now \((2n/p)=1\), and the identical root construction gives \(1\leq r<p\), \(2r^2=o(n)\), and a composite \(n-2r^2>p\).

Thus every sufficiently large \(n\) fails the property. Allowing \(x=0\) as well only makes the original property stronger and does not affect this disproof.

\paragraph{Attribution.}
Epure--Gica and Mollin--Williams give the more precise classification discussed on the problem page. The general Pollack-based approach is attributed in problem 1141 to the internal OpenAI model described in \url{https://arxiv.org/abs/2604.06609}; no new resolution or independent proof of Pollack's theorem is claimed.

\paragraph{Completion estimate.}
The infinitude question is fully answered negatively, conditional only on the stated established prime-character theorem and standard quadratic reciprocity and Bertrand's postulate.
\textbf{COMPLETION: 100\%}
